\documentclass[12pt, letterpaper]{article}

% --- Idioma y fuentes ---
\usepackage[utf8]{inputenc}
\usepackage[spanish, es-tabla]{babel}
\usepackage[T1]{fontenc}
\usepackage{lmodern}

% --- Geometría de página ---
\usepackage[top=2.5cm, bottom=2.5cm, left=2.5cm, right=2.5cm]{geometry}

% --- Matemáticas, tablas y gráficos ---
\usepackage{amsmath}
\usepackage{booktabs}
\usepackage{tabularx}
\usepackage{graphicx}
\usepackage{float}
\usepackage{fancyhdr}

% --- Colores (sin la opción xcdraw que generaba bloques negros) ---
\usepackage[table]{xcolor}

% --- Paleta en tonos pastel suaves ---
\definecolor{brandBlue}{RGB}{27, 79, 114}      % Azul elegante para títulos
\definecolor{darkText}{RGB}{44, 62, 80}        % Gris azulado oscuro para texto legible
\definecolor{pastelGreen}{RGB}{232, 248, 245}  % Verde menta pastel claro
\definecolor{pastelBlue}{RGB}{235, 245, 251}   % Azul pastel claro
\definecolor{pastelYellow}{RGB}{254, 249, 231} % Amarillo pastel suave
\definecolor{pastelGray}{RGB}{248, 249, 249}   % Fondo neutro grisáceo suave

% --- Cajas de notas estilizadas ---
\usepackage[most]{tcolorbox}
\newtcolorbox{cajaPastel}[2][]{
    boxrule=0.5pt, 
    arc=3pt, 
    colback=pastelBlue,
    colframe=brandBlue!60, 
    fonttitle=\bfseries\small\color{darkText},
    title={#2}, 
    #1
}

% --- Enlaces e hipervínculos ---
\usepackage{hyperref}
\hypersetup{
    colorlinks=true,
    linkcolor=brandBlue,
    urlcolor=brandBlue
}

% --- Encabezados y pies de página ---
\setlength{\headheight}{14.5pt}
\pagestyle{fancy}
\fancyhf{}
\lhead{\color{gray}\footnotesize Estructuras Computacionales --- Control de Pantalla RGB LCD}
\rhead{\color{gray}\footnotesize Tang Primer 20K}
\cfoot{\thepage}

% ===========================================================================
% INICIO DEL DOCUMENTO
% ===========================================================================
\begin{document}

% ===========================================================================
% PORTADA
% ===========================================================================
\begin{titlepage}
    \centering
    \vspace*{1.5cm}
    
    {\scshape\Large Universidad Nacional de Colombia \par}
    \vspace{0.3cm}
    {\large Sede Manizales \par}
    
    \vspace{2.5cm}
    \rule{\linewidth}{0.8pt} \\[0.4cm]
    {\Huge \bfseries\color{brandBlue} Estructuras Computacionales \par}
    \vspace{0.4cm}
    {\Large \bfseries\color{darkText} Implementación de Interfaz RGB LCD en FPGA \par}
    \rule{\linewidth}{0.8pt} \\[2cm]
    
    \vfill
    
    \begin{minipage}{0.8\textwidth}
        \centering
        \textbf{Presentado por:}\\[0.3em]
        {\large Marlyn Nathalia Mora Riascos}\\[0.2em]
        {\large Yeison}\\[1.5cm]
        
        \textbf{Docente:}\\[0.3em]
        {\large Prof. Gustavo Osorio}\\[1.5cm]
        
        \textbf{Fecha:}\\[0.3em]
        25 de septiembre de 2026
    \end{minipage}
    
    \vspace*{1.5cm}
\end{titlepage}

\newpage

% ===========================================================================
% 1. INTRODUCCIÓN Y PARADIGMAS DE DISEÑO
% ===========================================================================
\section{Introducción y Paradigmas de Diseño}

En el marco de la asignatura Estructuras Computacionales, el control de periféricos gráficos aborda dos niveles fundamentales de abstracción:

\begin{enumerate}
    \item \textbf{Capa de Hardware (Fase Actual -- Verilog):} Descripción a nivel de transferencia de registros (RTL) del controlador de video, generación de reloj de píxel (\textit{pixel clock}) mediante PLL y máquinas de sincronismo temporal (\textit{HSYNC}, \textit{VSYNC}, \textit{DE}).
    
    \item \textbf{Capa de Software (Fase Futura -- Ensamblador):} Programación en código máquina / Ensamblador para la arquitectura del procesador RISC-V (\texttt{Femto\_Risc-V\_SoC}), permitiendo escribir caracteres y gráficos mediante Entrada/Salida Mapeada en Memoria (MMIO).
\end{enumerate}

\begin{table}[H]
\centering
\small
\renewcommand{\arraystretch}{1.3}
\begin{tabularx}{\textwidth}{>{\columncolor{pastelGreen}}p{4.2cm} >{\columncolor{pastelGray}}X}
\toprule
\textbf{\color{brandBlue}Paradigma / Lenguaje} & \textbf{\color{brandBlue}Rol en la Arquitectura del Sistema} \\
\midrule
\textbf{Verilog (HDL / RTL)} & Modela la circuitería física en la FPGA. Construye el controlador de pantalla que genera los pulsos de barrido y maneja el bus físico de 40 pines. \\
\midrule
\textbf{Ensamblador (Assembly)} & Instrucciones de bajo nivel ejecutadas por la CPU. Controla el flujo del programa y escribe datos en los registros del controlador de video. \\
\bottomrule
\end{tabularx}
\caption{Diferenciación de capas en el diseño del sistema embebido.}
\label{tab:paradigmas}
\end{table}

\subsection{Plataforma de Hardware}

Para validar la capa física se emplearon los siguientes componentes:

\begin{table}[H]
\centering
\small
\renewcommand{\arraystretch}{1.3}
\begin{tabularx}{\textwidth}{>{\columncolor{pastelBlue}}p{4.2cm} >{\columncolor{pastelGray}}X}
\toprule
\textbf{\color{brandBlue}Elemento} & \textbf{\color{brandBlue}Especificación y Referencia Exacta} \\
\midrule
\textbf{Tarjeta FPGA} & Sipeed Tang Primer 20K (Gowin GW2A-LV18PG256C8/I7). \\
\textbf{Tarjeta Base (Dock)} & Carrier Dock con conector FPC de 40 pines y regulador elevador (\textit{boost}) para backlight. \\
\textbf{Módulo de Pantalla} & Panel TFT LCD de 7.0'' \textbf{SH070JPB60} (Lote: SH070JPB60-07847L-88). \\
\textbf{Cinta Flexible (FPC)} & 40 pines, paso de $0.5\text{ mm}$: \textbf{FPC-070JP27-40} (Marcado: 070JPB148-40 BXF-11-HX-B). \\
\textbf{Resolución y Timings} & $800 \times 480$ píxeles a $60\text{ Hz}$ (reloj de píxel de $40\text{ MHz}$). \\
\bottomrule
\end{tabularx}
\caption{Plataforma de hardware.}
\label{tab:hardware}
\end{table}

% ===========================================================================
% 2. ORIGEN DEL CÓDIGO Y ESTRUCTURA DE MÓDULOS
% ===========================================================================
\section{Origen del Código y Función de Módulos}

Para el desarrollo del controlador se tomó como base de referencia el repositorio oficial de Sipeed para la Tang Primer 20K:

\begin{center}
    \small 
    \url{https://github.com/sipeed/TangPrimer-20K-example} \\
    $\rightarrow$ \textbf{Ruta:} \texttt{RGB\_lcd/rgb\_lcd\_5inch\_colorbar/}
\end{center}

Aunque la referencia del fabricante indica 5 pulgadas, la pantalla (\textbf{SH070JPB60}) implementa la misma resolución nativa ($800 \times 480$), idénticos parámetros de refresco a $60\text{ Hz}$ y el mismo conector FPC de 40 pines.

\begin{table}[H]
\centering
\small
\renewcommand{\arraystretch}{1.3}
\begin{tabularx}{\textwidth}{>{\columncolor{pastelYellow}}p{3.4cm} >{\columncolor{pastelGray}}X}
\toprule
\textbf{\color{brandBlue}Archivo / Módulo} & \textbf{\color{brandBlue}Función Específica en la Arquitectura} \\
\midrule
\textbf{\texttt{top.v}} & Módulo integrador (\textit{Top-level}). Conecta el oscilador, instancia el PLL, el controlador de barrido y mapea los buses de color a las salidas físicas. \\
\midrule
\textbf{\texttt{lcd\_ctrl.v}} & Controlador de barrido de video. Cuenta ciclos de reloj para delimitar las líneas horizontales y cuadros verticales, generando las señales \texttt{HSYNC}, \texttt{VSYNC} y \texttt{DE} (\textit{Data Enable}). \\
\midrule
\textbf{\texttt{lcd\_data.v}} & Generador de patrones gráficos. Recibe la coordenada $(X, Y)$ del controlador y define qué valor RGB (rojo, verde, azul, blanco, etc.) transmitir en cada píxel. \\
\midrule
\textbf{\texttt{pll\_40m.v}} & Primitiva de reloj \texttt{rPLL} de Gowin. Eleva la frecuencia del oscilador de la tarjeta ($27\text{ MHz}$) al reloj de píxel requerido por la pantalla ($40\text{ MHz}$). \\
\midrule
\textbf{\texttt{rgb\_lcd.cst}} & Archivo de restricciones físicas (\textit{Physical Constraints}). Asocia cada señal lógica del código a los pines BGA físicos del conector FPC de 40 pines. \\
\midrule
\textbf{\texttt{Makefile}} & Script de automatización. Orquesta las etapas de síntesis, ruteo y flasheo mediante comandos reproducibles dentro del contenedor Docker. \\
\bottomrule
\end{tabularx}
\caption{Descripción funcional de los módulos del proyecto.}
\label{tab:modulos}
\end{table}


% ===========================================================================
% 3. FLUJO DE SÍNTESIS Y COMPILACIÓN (TOOLCHAIN)
% ===========================================================================
\section{Flujo de Síntesis y Compilación}

Para evitar conflictos de dependencias, todo el proceso de compilación se encapsuló en un contenedor Docker (\texttt{fpga-toolchain}), articulando cuatro herramientas abiertas especializadas:

\begin{enumerate}
    \item \textbf{Síntesis Lógica (\texttt{Yosys}):} Analiza la descripción en Verilog y traduce el diseño a una red de compuertas lógicas y tablas de búsqueda (LUTs) específicas de Gowin (\texttt{top\_synth.json}).
    
    \item \textbf{Ubicación y Ruteo (\texttt{nextpnr-himbaechel}):} Asigna las compuertas a las coordenadas físicas del chip GW2A-18 y rutea las pistas de interconexión respetando las restricciones de pines del archivo \texttt{rgb\_lcd.cst} (\texttt{top\_pnr.json}).
    
    \item \textbf{Empaquetado (\texttt{gowin\_pack}):} Convierte la base de datos de ruteo al formato binario nativo (\textit{Bitstream}) de la FPGA: \textbf{\texttt{top.fs}}.
    
    \item \textbf{Descarga a Hardware (\texttt{openFPGALoader}):} Envía el archivo \texttt{top.fs} a través del puerto USB hacia la memoria SRAM de la FPGA utilizando el protocolo JTAG mediante el chip FT2232.
\end{enumerate}

\begin{table}[H]
\centering
\small
\renewcommand{\arraystretch}{1.3}
\begin{tabularx}{\textwidth}{>{\columncolor{pastelBlue}}p{4.2cm} >{\columncolor{pastelGray}}X}
\toprule
\textbf{\color{brandBlue}Comando (\texttt{Makefile})} & \textbf{\color{brandBlue}Operación Ejecutada} \\
\midrule
\texttt{make detect} & Verifica la comunicación JTAG con el chip Gowin GW2A-18 vía USB. \\
\midrule
\texttt{make synth} & Ejecuta la síntesis RTL con Yosys y genera el archivo JSON de compuertas. \\
\midrule
\texttt{make bitstream} & Corre el ruteo (\textit{PnR}) con nextpnr y empaqueta el binario final \texttt{top.fs}. \\
\midrule
\texttt{make flash} & Programa la SRAM de la FPGA mediante openFPGALoader para prueba inmediata. \\
\bottomrule
\end{tabularx}
\caption{Comandos del flujo de trabajo automatizado.}
\label{tab:toolchain}
\end{table}


% ===========================================================================
% 4. REPOSITORIO Y CÓDIGO FUENTE
% ===========================================================================
\section{Repositorio y Código Fuente}

El código fuente completo del proyecto, incluyendo los módulos en Verilog (\texttt{top.v}, \texttt{lcd\_ctrl.v}, \texttt{lcd\_data.v}, \texttt{pll\_40m.v}), el archivo de restricciones físicas (\texttt{rgb\_lcd.cst}) y el \texttt{Makefile} de compilación, se encuentra alojado y versionado en el siguiente repositorio:

\begin{tcolorbox}[
    colback=pastelBlue, 
    colframe=brandBlue!60, 
    arc=3pt,
    boxrule=0.6pt,
    title=\textbf{Acceso al Repositorio del Proyecto}
]
    \centering
    \vspace{0.2cm}
    \textbf{Enlace al código fuente completo:}\\[0.4em]
    % Reemplaza esta URL por el enlace real de tu repositorio:
    \href{https://github.com/tu-usuario/tu-repositorio}{\texttt{https://github.com/tu-usuario/tu-repositorio}}
    \vspace{0.2cm}
\end{tcolorbox}

% ===========================================================================
% 5. RESULTADOS Y VALIDACIÓN EXPERIMENTAL
% ===========================================================================
\section{Resultados y Validación Experimental}

Se cargó el \textit{bitstream} \texttt{top.fs} en la memoria SRAM de la FPGA mediante el comando \texttt{make flash}, logrando el encendido del panel TFT LCD y la visualización nítida y estable de las franjas de colores.

\begin{figure}[H]
    \centering
    % Ajusta width al tamaño deseado (ej. 0.75\textwidth o 0.7\textwidth)
    \includegraphics[width=0.72\textwidth]{"WhatsApp Image 2026-09-25 at 7.37.08 PM.jpeg"}
    \caption{Validación experimental del controlador RGB LCD en funcionamiento sobre la tarjeta Tang Primer 20K con el panel TFT.}
    \label{fig:resultado}
\end{figure}

% ===========================================================================
% 6. CONCLUSIÓN
% ===========================================================================
\section{Conclusión}

Se integró y validó con éxito el controlador de video RGB sobre la FPGA Sipeed Tang Primer 20K para la pantalla SH070JPB60 a una resolución de $800 \times 480$ píxeles. La arquitectura implementada en Verilog proporciona una base determinística en hardware sobre la cual se desarrollará posteriormente el control por software en ensamblador mediante el procesador RISC-V del curso.

\end{document}